\documentclass[11pt,a4paper]{article}
\usepackage[T1]{fontenc}
\usepackage[margin=25mm]{geometry}
\usepackage{amsmath,amssymb,amsthm}
\usepackage{microtype}

\newtheorem{conjecture}{Conjecture}
\pagestyle{empty}
\setlength{\parindent}{0pt}
\setlength{\parskip}{0.45em}

\title{The Geodesic Deletion Conjecture}
\author{Luc Gillibert, GPT 6 Astra}
\date{17/11/2026}

\begin{document}
\maketitle
\vspace{-2em}

All graphs in this note are finite, simple, and undirected. A graph is
\emph{vertex-transitive} if, for any two vertices, an automorphism maps
one to the other. A \emph{geodesic} is a shortest path between its
endpoints. We regard a graph consisting of a single vertex as connected.

\begin{conjecture}[Geodesic Deletion]
Let $G$ be a connected vertex-transitive graph with at least three
vertices. For every geodesic $P$ in $G$, the graph
\[
  \boxed{G-V(P)\text{ is connected}.}
\]
\end{conjecture}

Here $G-V(P)$ denotes the induced subgraph on
$V(G)\setminus V(P)$: all vertices of $P$, including both endpoints,
are deleted, together with their incident edges. Thus the conjecture
asserts that \emph{every} shortest route can be removed without
disconnecting the remaining vertices. A path with this property is
called \emph{nonseparating}.

\paragraph{Elementary examples.}
For a cycle $C_n$ with $n\geq 3$, a geodesic occupies a consecutive arc
of length at most $\lfloor n/2\rfloor$. Deleting its vertices leaves a
nonempty path, possibly a single vertex. For a complete graph $K_n$
with $n\geq 3$, a geodesic has at most two vertices, so its deletion
leaves a nonempty complete graph.

\paragraph{Edge-transitivity is insufficient.}
The complete bipartite graph $K_{2,3}$ is edge-transitive. Write its
parts as $\{a,b\}$ and $\{x,y,z\}$. The path $a,x,b$ is a geodesic,
since $a$ and $b$ are nonadjacent and have distance two. Deleting its
vertices leaves the two isolated vertices $y$ and $z$. Hence the
vertex-transitivity hypothesis cannot simply be replaced by
edge-transitivity.

\paragraph{The shortest-path condition is essential.}
In the cube $Q_3$, whose vertices are binary strings of length three
and whose edges join strings differing in exactly one coordinate,
consider the induced path
\[
  000,\ 001,\ 011,\ 111,\ 110.
\]
Its deletion leaves the isolated vertex $010$ and the separate edge
$100\text{--}101$. This path has length four, whereas its endpoints
have distance two. Thus even in a vertex-transitive graph, deleting
an induced path need not leave a connected graph.

\paragraph{Motivation.}
The conjecture links a metric condition on a single path to a global
connectivity property of its complement. Vertex-transitivity alone
does not make all geodesics equivalent under automorphisms, so the
universal quantifier over shortest paths is a substantive requirement.


\end{document}
